% Agentic Coding Overview, 30 September 2026
% Build with: latexmk main.tex
\documentclass[aspectratio=169]{beamer}
\usepackage{beamerthemeACIN}
\usepackage{booktabs}
\usepackage{array}
\newcolumntype{P}[1]{>{\raggedright\arraybackslash}p{#1}}
\usepackage{hyperref}
\usetikzlibrary{positioning}
\hypersetup{colorlinks=true,urlcolor=TUblue}

\title{Agentic Coding Overview}
\subtitle{Models, harnesses, access, routing, and benchmarks}
\author{Cédric Jung \and Jozef Janus}
\date{30 September 2026}

\newcommand{\src}[2]{\par\vfill{\fontsize{7}{8}\selectfont\color{ACINgray}Source: \href{#1}{#2}}}
\newcommand{\tagcurrent}{{\color{TUblue}\bfseries Listed}}
\newcommand{\tagcandidate}{{\color{ACINred}\bfseries Candidate}}

\begin{document}

{% Title page has its own logo treatment.
\setbeamertemplate{footline}{}
\begin{frame}[plain]
  \titlepage
\end{frame}
}

\begin{frame}{AQUEDUCT: models listed today}{Snapshot from the public overview, 30 September 2026}
\small
\renewcommand{\arraystretch}{1.35}
\begin{tabular}{@{}P{4.5cm}P{2.4cm}P{4.7cm}@{}}
\toprule
\textbf{Model} & \textbf{Context} & \textbf{Deployment note} \\
\midrule
DeepSeek V4 Flash 284B & 393k tokens & RTX PRO 6000 Blackwell migration \\
Qwen 3.6 35B & 262k tokens & 4 replicas, one A40 each \\
GLM 5.2 744B Preview & 262k tokens & 8 AMD MI300X GPUs \\
\bottomrule
\end{tabular}

\vspace{4mm}
\begin{block}{Transition}
Qwen 3.5 397B reached its announced retirement date on 30 September. DeepSeek V4 Flash is its planned replacement on the Blackwell infrastructure.
\end{block}
\src{https://datalab.tuwien.ac.at/aiml/aqueduct/}{TU Wien AQUEDUCT overview and release notes}
\end{frame}

\begin{frame}{AQUEDUCT: what may come next}{Release-note candidates are not confirmed deployments}
\footnotesize
\renewcommand{\arraystretch}{1.18}
\begin{tabular}{@{}P{4.5cm}P{4.0cm}P{3.1cm}@{}}
\toprule
\textbf{Candidate model} & \textbf{Discussed hardware} & \textbf{Status} \\
\midrule
Qwen3.8-27B & freed H200 node; possibly 4 instances & \tagcandidate \\
DeepSeek V4 Flash Vision Exp & RTX PRO 6000 Blackwell & \tagcandidate \\
GLM-5.3-Flash & RTX PRO 6000 Blackwell & \tagcandidate \\
DeepSeek V4.1 Flash / GLM-5.3 & AMD MI300X & \tagcandidate \\
\bottomrule
\end{tabular}

\vspace{2mm}
\begin{exampleblock}{Before use}
Check the live \texttt{/v1/models} endpoint. A release-note mention does not make a model available through AQUEDUCT.
\end{exampleblock}
\src{https://datalab.tuwien.ac.at/aiml/aqueduct/release-notes/}{TU Wien AQUEDUCT release notes}
\end{frame}

\begin{frame}{Hardware behind the AQUEDUCT models}{Memory and replica counts shape what can be served}
\small
\renewcommand{\arraystretch}{1.38}
\begin{tabular}{@{}P{3.55cm}P{2.05cm}P{6.0cm}@{}}
\toprule
\textbf{Accelerator} & \textbf{Memory / GPU} & \textbf{AQUEDUCT use} \\
\midrule
NVIDIA A40 & 48 GB & Qwen 3.6 35B: 4 independent one-GPU replicas \\
RTX PRO 6000 Blackwell & 96 GB & DeepSeek V4 Flash replaces the retiring Qwen 3.5 deployment \\
NVIDIA H200 & 141 GB & Initially served DeepSeek; freed for a future deployment \\
AMD MI300X & 192 GB & GLM 5.2 744B Preview: 8 GPUs in one replica \\
\bottomrule
\end{tabular}

\vspace{2mm}
{\color{ACINgray}\footnotesize Hardware mapping reflects the public overview plus migration notes; verify the live deployment before relying on it.}
\src{https://datalab.tuwien.ac.at/aiml/aqueduct/}{TU Wien AQUEDUCT overview; \href{https://datalab.tuwien.ac.at/aiml/aqueduct/release-notes/}{release notes}}
\end{frame}

\begin{frame}{Student access to AQUEDUCT}{Access requires a request and allowlisting}
\begin{enumerate}\itemsep8pt
  \item \textbf{Request access} through the TU Wien AQUEDUCT access process.
  \item \textbf{Wait for approval / allowlisting} before using the service.
  \item \textbf{Use the issued credentials} in an API-compatible client or coding harness.
  \item \textbf{Check available models} through the live \texttt{/v1/models} endpoint.
\end{enumerate}
\vspace{3mm}
\begin{block}{Practical point}
The public model catalog describes the service; it does not grant an account automatically.
\end{block}
\src{https://datalab.tuwien.ac.at/aiml/aqueduct/}{TU Wien AQUEDUCT service page; access requirement from presentation brief}
\end{frame}

\begin{frame}{Model and harness: two parts of a coding agent}
\begin{columns}[T,onlytextwidth]
\begin{column}{.47\textwidth}
  \begin{block}{Model: decides}
  \begin{itemize}\itemsep5pt
    \item Understands the task and code
    \item Proposes edits and tool calls
    \item Interprets tool results
  \end{itemize}
  \end{block}
  \centering\small Examples: GPT, Claude, Qwen, DeepSeek
\end{column}
\hfill
\begin{column}{.47\textwidth}
  \begin{exampleblock}{Harness: executes the loop}
  \begin{itemize}\itemsep5pt
    \item Builds context and exposes tools
    \item Reads files, edits code, runs tests
    \item Manages state and permissions
  \end{itemize}
  \end{exampleblock}
  \centering\small Examples: Codex, Claude Code, OpenCode
\end{column}
\end{columns}
\vspace{6mm}
\centering\large User task $\rightarrow$ harness $\leftrightarrow$ model $\rightarrow$ tested change
\src{https://openai.com/index/unrolling-the-codex-agent-loop/}{OpenAI, \emph{Unrolling the Codex agent loop}}
\end{frame}

\begin{frame}{Post-training can include the agent loop}{Codex as a concrete example}
\small
\begin{center}
\begin{tikzpicture}[>=stealth, every node/.style={font=\small}, node distance=1.4cm]
  \node[draw=ACINred,fill=ACINred!6,rounded corners,minimum width=2.2cm,minimum height=.9cm] (task) {Coding task};
  \node[draw=TUblue,fill=TUblue!6,rounded corners,minimum width=2.5cm,minimum height=.9cm,right=of task] (agent) {Model + tools};
  \node[draw=ACINred,fill=ACINred!6,rounded corners,minimum width=2.6cm,minimum height=.9cm,right=of agent] (test) {Test feedback};
  \draw[->,thick] (task) -- (agent);
  \draw[->,thick] (agent) -- (test);
  \draw[->,thick] (test.south) .. controls +(0,-.7) and +(0,-.7) .. node[below]{next action / training signal} (agent.south);
\end{tikzpicture}
\end{center}
\vspace{3mm}
\begin{itemize}\itemsep6pt
  \item Coding post-training can use trajectories: inspect, edit, run tests, observe, revise.
  \item OpenAI describes codex-1 and GPT-5-Codex as trained or optimized for agentic coding in tool-using environments.
  \item A benchmark therefore measures the \textbf{model, harness, tools, and reasoning budget} together.
\end{itemize}
\src{https://openai.com/index/gpt-5-system-card-addendum-gpt-5-codex/}{OpenAI, GPT-5-Codex system-card addendum}
\end{frame}

\begin{frame}{Why LLM routing platforms exist}{One integration, several model and provider choices}
\begin{center}
\begin{tikzpicture}[>=stealth, every node/.style={font=\small}]
  \node[draw=ACINgray,rounded corners,minimum width=2cm,minimum height=.9cm] (app) at (0,0) {Application};
  \node[draw=ACINred,fill=ACINred!6,rounded corners,minimum width=2cm,minimum height=.9cm] (router) at (4,0) {Router};
  \node[draw=TUblue,rounded corners,minimum width=2.4cm] (p1) at (8,1.2) {Provider A};
  \node[draw=TUblue,rounded corners,minimum width=2.4cm] (p2) at (8,0) {Provider B};
  \node[draw=TUblue,rounded corners,minimum width=2.4cm] (p3) at (8,-1.2) {Provider C};
  \draw[->,thick] (app) -- node[above]{one API} (router);
  \draw[->] (router) -- (p1); \draw[->] (router) -- (p2); \draw[->] (router) -- (p3);
\end{tikzpicture}
\end{center}
\vspace{3mm}
\begin{itemize}\itemsep4pt
  \item \textbf{Model} = which neural network responds.
  \item \textbf{Provider} = who runs it; providers vary in price, speed, region, and retention rules.
  \item The router can select a provider, apply constraints, and fail over when needed.
\end{itemize}
\src{https://openrouter.ai/docs/quickstart}{OpenRouter quickstart; \href{https://www.eurouter.ai/docs/concepts/routing}{EUrouter routing docs}}
\end{frame}

\begin{frame}{OpenRouter}{Global catalog and provider routing}
\begin{columns}[T,onlytextwidth]
\begin{column}{.48\textwidth}
\begin{block}{What it offers}
\begin{itemize}\itemsep5pt
  \item One API across many models
  \item Provider ordering and filtering
  \item Price or throughput priority
  \item Provider and model fallback
\end{itemize}
\end{block}
\end{column}
\hfill
\begin{column}{.48\textwidth}
\begin{exampleblock}{Routing choices}
\begin{itemize}\itemsep5pt
  \item Cheapest eligible endpoint
  \item Higher throughput endpoint
  \item Specific provider or region
  \item EU in-region option for eligible enterprise use
\end{itemize}
\end{exampleblock}
\end{column}
\end{columns}
\vspace{5mm}
\centering\small \textbf{Check the chosen provider and region} before sending sensitive code or data.
\src{https://openrouter.ai/docs/quickstart}{OpenRouter quickstart; \href{https://openrouter.ai/docs/guides/get-started/sovereign-ai}{sovereign AI docs}}
\end{frame}

\begin{frame}{European routing and hosting options}{Different ways to constrain where inference runs}
\small
\renewcommand{\arraystretch}{1.4}
\begin{tabular}{@{}P{2.7cm}P{4.6cm}P{4.8cm}@{}}
\toprule
\textbf{Service} & \textbf{Role} & \textbf{Regional control} \\
\midrule
EUrouter & Multi-provider model gateway & EU processing by design; provider filters \\
Eden AI & Multi-provider AI gateway & Dedicated European endpoint \\
Scaleway & Inference provider with many models & European-hosted serving; not a multi-provider router \\
\bottomrule
\end{tabular}
\vspace{5mm}
\begin{block}{Selection question}
Specify processing location, provider ownership, and retention requirements separately. ``EU-hosted'' and ``EU-owned'' are different constraints.
\end{block}
\src{https://www.eurouter.ai/docs}{EUrouter docs; \href{https://www.edenai.co/eu}{Eden AI EU}; \href{https://www.scaleway.com/en/docs/generative-apis/faq/}{Scaleway FAQ}}
\end{frame}

\begin{frame}{Subscription versus API pricing}{Illustrative September 2026 snapshot}
\begin{columns}[T,onlytextwidth]
\begin{column}{.48\textwidth}
  \begin{block}{Subscription: fixed monthly fee}
  \textbf{ChatGPT Plus: \$20/month}\par
  Codex usage is included within plan limits. Good fit for interactive individual use.
  \end{block}
\end{column}
\hfill
\begin{column}{.48\textwidth}
  \begin{exampleblock}{API: metered tokens}
  \textbf{GPT-5.6 Sol: \$4/M input + \$20/M output}\par
  Example: 5M input + 1M output $=\$40$. Good fit for programmable workflows.
  \end{exampleblock}
\end{column}
\end{columns}
\vspace{5mm}
\small The subscription is not an API credit balance: it has product and usage limits. Caching and workload shape API cost. Recheck rates before purchasing.
\src{https://help.openai.com/en/articles/6950777-what-is-chatgpt-plus}{ChatGPT Plus; \href{https://developers.openai.com/api/docs/models/gpt-5.6-sol}{GPT-5.6 Sol API pricing}}
\end{frame}

\begin{frame}{Which benchmarks answer which questions?}{Keep versions, metrics, and evaluators visible}
\footnotesize
\renewcommand{\arraystretch}{1.2}
\begin{tabular}{@{}P{3.3cm}P{3.7cm}P{5.0cm}@{}}
\toprule
\textbf{Benchmark} & \textbf{Measure} & \textbf{Use / limitation} \\
\midrule
GDPval-AA v2.1 & Elo & Written knowledge work; proxy, not academic paper writing \\
DeepSWE v1.1 & \% tasks resolved & Long-horizon software engineering; harness varies \\
Terminal-Bench 4.0 & \% tasks passed & Agentic terminal work; common evaluator for chart \\
Humanity's Last Exam & pass@1 \% & Broad expert reasoning; not science-only \\
\bottomrule
\end{tabular}
\vspace{2mm}
\begin{block}{Reading rule}
Compare exact model revisions under the same benchmark version. Missing results mean \emph{no public exact-model score located}, not zero.
\end{block}
\src{https://www.tbench.ai/}{Terminal-Bench; \href{https://deepswe.datacurve.ai/}{DeepSWE}; \href{https://artificialanalysis.ai/evaluations/gdpval-aa}{GDPval-AA}}
\end{frame}

\begin{frame}{Agentic coding result: Terminal-Bench 4.0}{Approximate task pass rate, Artificial Analysis setup}
\centering
\begin{tikzpicture}[x=1cm,y=.52cm,every node/.style={font=\fontsize{8}{9}\selectfont}]
  \draw[ACINlightgray] (.15,-.25) -- (6.8,-.25);
  \foreach \t in {0,20,40,60} {
    \draw[ACINlightgray] ({.15+.105*\t},-.25) -- ({.15+.105*\t},7.05);
    \node[anchor=north,ACINgray] at ({.15+.105*\t},-.3) {\t\%};
  }
  \foreach \label/\score/\y/\col in {{DeepSeek V4 Flash 0731}/12/6/ACINred,{DeepSeek V4 Flash Vision}/12/5/ACINred,{GLM-5.3-Flash}/33/4/ACINred,{GPT-5.6 Sol}/40/3/TUblue,{Claude Opus 5.5}/59.6/2/TUblue,{Grok 4.7}/26/1/TUblue,{Gemini 3.8 Flash}/20/0/TUblue} {
    \node[anchor=east] at (0,\y) {\label};
    \fill[fill=\col] (.15,\y-.26) rectangle ({.15+.105*\score},\y+.26);
    \node[anchor=west] at ({.25+.105*\score},\y) {\score\%};
  }
\end{tikzpicture}
\par\vspace{1mm}
{\fontsize{8}{9}\selectfont\color{ACINgray}Red: open-weight comparison models; blue: proprietary frontier models. Qwen3.6-35B-A3B has no exact v4.0 result. Scores are approximate and depend on harness and reasoning settings.}
\src{https://artificialanalysis.ai/evaluations/terminalbench-4-0}{Artificial Analysis, Terminal-Bench 4.0; model-name mapping in supplied benchmark notes}
\end{frame}

\begin{frame}{Sources and further reading}{Service, agents, and routing}
\footnotesize
\begin{itemize}\itemsep5pt
  \item \href{https://datalab.tuwien.ac.at/aiml/aqueduct/}{TU Wien AQUEDUCT overview} and \href{https://datalab.tuwien.ac.at/aiml/aqueduct/release-notes/}{release notes}
  \item \href{https://openai.com/index/unrolling-the-codex-agent-loop/}{OpenAI: Unrolling the Codex agent loop}
  \item \href{https://openai.com/index/gpt-5-system-card-addendum-gpt-5-codex/}{OpenAI: GPT-5-Codex system-card addendum}
  \item \href{https://openrouter.ai/docs/quickstart}{OpenRouter quickstart} and \href{https://openrouter.ai/docs/guides/get-started/sovereign-ai}{regional routing}
  \item \href{https://www.eurouter.ai/docs}{EUrouter documentation}, \href{https://www.edenai.co/eu}{Eden AI EU}, \href{https://www.scaleway.com/en/docs/generative-apis/faq/}{Scaleway FAQ}
  \item \href{https://help.openai.com/en/articles/6950777-what-is-chatgpt-plus}{ChatGPT Plus} and \href{https://developers.openai.com/api/docs/models/gpt-5.6-sol}{GPT-5.6 Sol API pricing}
\end{itemize}
\vspace{4mm}
\color{ACINgray}\scriptsize Deck compiled from the repository's \texttt{documentation/} research notes, dated 30 September 2026. Links are clickable in the PDF.
\end{frame}

\begin{frame}{Benchmark sources}{Exact versions and configurations matter}
\footnotesize
\begin{itemize}\itemsep6pt
  \item \href{https://artificialanalysis.ai/evaluations/gdpval-aa}{Artificial Analysis: GDPval-AA v2.1}
  \item \href{https://deepswe.datacurve.ai/}{DataCurve: DeepSWE v1.1}
  \item \href{https://artificialanalysis.ai/evaluations/terminalbench-4-0}{Artificial Analysis: Terminal-Bench 4.0}
  \item \href{https://artificialanalysis.ai/evaluations/humanitys-last-exam}{Artificial Analysis: Humanity's Last Exam}
  \item \href{https://www.tbench.ai/}{Terminal-Bench project}
\end{itemize}
\vspace{4mm}
\begin{block}{Interpretation}
Agentic scores reflect model, harness, tools, and reasoning budget. Benchmark versions and evaluators must match for a fair comparison.
\end{block}
\end{frame}

\end{document}
